\documentclass[12pt,letterpaper]{article}
\usepackage[utf8]{inputenc}
\usepackage[T1]{fontenc}
\usepackage[spanish]{babel}
\usepackage[margin=2.5cm]{geometry}
\usepackage{graphicx}
\usepackage{booktabs}
\usepackage{tikz}
\usetikzlibrary{shapes.geometric,arrows.meta,positioning,fit,backgrounds}
\usepackage{hyperref}
\usepackage{url}
\usepackage{xcolor}
\usepackage{setspace}
\usepackage{titlesec}

\hypersetup{colorlinks=true, urlcolor=blue, linkcolor=blue, citecolor=blue}
\sloppy
\onehalfspacing

\titleformat{\section}{\large\bfseries}{\thesection.}{0.5em}{}
\titleformat{\subsection}{\normalsize\bfseries}{\thesubsection}{0.5em}{}

\begin{document}

% ============================================================
% PORTADA (Seccion 1 de la estructura obligatoria - no cuenta
% para el minimo de 5 paginas de contenido)
% ============================================================
\begin{titlepage}
\centering
\vspace*{1.5cm}

{\large UNIVERSIDAD T\'ECNICA ESTATAL DE QUEVEDO\par}
{\normalsize Facultad de Ciencias de la Computaci\'on\par}
{\normalsize Carrera de Ingenier\'ia de Software\par}
\vspace{0.3cm}
{\normalsize Aplicaciones Distribuidas (ISR-701) --- Unidad 3\par}

\vspace{2cm}
{\Large\bfseries An\'alisis Cr\'itico del Modelo de Consistencia \\
y Replicaci\'on en el Sistema de Control de \\
Laboratorios Inform\'aticos (SCLI)\par}

\vspace{0.5cm}
{\normalsize CockroachDB, Consenso Raft y Coordinaci\'on entre Microservicios\par}

\vspace{2cm}
{\normalsize\bfseries C\'odigo: TA-IND-02\par}
{\normalsize Trabajo Aut\'onomo Individual --- Sumativo\par}

\vspace{1.5cm}
{\normalsize\bfseries Estudiante:\par}
{\normalsize Harold Nicol\'as Vinueza S\'anchez\par}

\vspace{0.8cm}
{\normalsize\bfseries Equipo PFC:\par}
{\normalsize FUVV --- Sistema de Control de Laboratorios Inform\'aticos (SCLI)\par}

\vspace{0.8cm}
{\normalsize\bfseries Docente:\par}
{\normalsize Ph.D. Gleiston C. Guerrero-Ulloa\par}

\vspace{1.5cm}
{\normalsize Per\'iodo Acad\'emico 2026--2027\par}
{\normalsize Quevedo, Ecuador --- Julio de 2026\par}

\vfill
\end{titlepage}

% ============================================================
% RESUMEN (Seccion 2)
% ============================================================
\section*{Resumen}
Este informe analiza cr\'iticamente el modelo de consistencia y replicaci\'on
implementado en el Sistema de Control de Laboratorios Inform\'aticos (SCLI),
proyecto fin de curso del equipo FUVV. SCLI adopta una arquitectura de
microservicios con base de datos dedicada por servicio sobre CockroachDB,
lo que produce dos capas de consistencia distintas: una capa
\emph{intra-servicio} fuerte, garantizada por el consenso Raft de
CockroachDB con aislamiento serializable, y una capa \emph{inter-servicio}
sin garant\'ias formales, sostenida \'unicamente por llamadas REST
s\'incronas. Se caracteriza el modelo real implementado, se eval\'uan dos
alternativas de consistencia y replicaci\'on frente al modelo actual
mediante criterios t\'ecnicos expl\'icitos, y se responde de forma razonada
si el modelo elegido es \'optimo para el dominio de reserva y monitoreo de
laboratorios. Se concluye que la fortaleza de consistencia intra-servicio
est\'a bien justificada, pero la ausencia de un mecanismo expl\'icito de
coordinaci\'on entre servicios constituye una limitaci\'on real que debe
resolverse antes de completar la integraci\'on del m\'odulo de reservas.

\vspace{0.5cm}

% ============================================================
% INTRODUCCION (Seccion 3)
% ============================================================
\section{Introducci\'on}

El Sistema de Control de Laboratorios Inform\'aticos (SCLI) es el proyecto
fin de curso (PFC) del equipo FUVV, correspondiente al dominio de referencia
\emph{Laboratorios Inform\'aticos} definido para el per\'iodo. El sistema
migra un monolito de gesti\'on de laboratorios hacia una arquitectura de
microservicios (\texttt{auth-service}, \texttt{usuarios-service},
\texttt{academico-laboratorios-service} y
\texttt{reservas-solicitudes-service}), cada uno con su propia base de
datos CockroachDB, orquestados con Docker Compose y comunicados mediante
REST interno protegido con una clave compartida.

Esta arquitectura de \emph{base de datos por servicio} plantea una pregunta
que no es meramente acad\'emica sino directamente aplicable al estado
actual del c\'odigo del equipo: \textbf{\textquestiondown es el modelo de
consistencia elegido el \'optimo para el dominio del sistema?} El presente
informe responde a esa pregunta caracterizando el modelo realmente
implementado --- tanto dentro de cada servicio como entre servicios ---
contrast\'andolo con al menos dos alternativas de dise\~no y discutiendo sus
compromisos a la luz del teorema CAP y del modelo PACELC.

% ============================================================
% MARCO TEORICO (Seccion 4)
% ============================================================
\section{Marco Te\'orico}

\subsection{Espectro de modelos de consistencia}

La consistencia de r\'eplicas no es una propiedad binaria sino un espectro
ordenado por fuerza sem\'antica \cite{viotti2016}. En el extremo m\'as
fuerte, la \textbf{linealizabilidad} (o consistencia fuerte) garantiza que
toda operaci\'on parece ejecutarse instant\'aneamente en un punto \'unico
entre su inicio y su fin, de modo que cualquier lector observa siempre el
\'ultimo valor escrito. La \textbf{consistencia secuencial} relaja el
requisito temporal pero preserva un orden global \'unico compartido por
todos los procesos. La \textbf{consistencia causal} solo exige preservar el
orden entre operaciones causalmente relacionadas, permitiendo reordenar
operaciones concurrentes independientes. La \textbf{consistencia de lectura
de los propios cambios} (\emph{read-your-writes}) garantiza \'unicamente
que un proceso vea sus propias escrituras previas. En el extremo m\'as
d\'ebil, la \textbf{consistencia eventual} solo promete que, en ausencia de
nuevas escrituras, todas las r\'eplicas converger\'an eventualmente al
mismo estado. Un modelo m\'as fuerte impone mayor coordinaci\'on entre
nodos --- y por tanto mayor latencia --- mientras que uno m\'as d\'ebil
admite mayor disponibilidad a costa de exponer estados transitoriamente
divergentes a la aplicaci\'on \cite{kleppmann2017}.

\subsection{Teoremas de compromiso: CAP y PACELC}

El teorema CAP, formalizado por Gilbert y Lynch a partir de la conjetura de
Brewer, establece que ante una partici\'on de red un sistema distribuido no
puede garantizar simult\'aneamente consistencia y disponibilidad
\cite{gilbert2002}. Dado que las particiones son inevitables en cualquier
red real, la decisi\'on de dise\~no efectiva es CP (priorizar consistencia)
o AP (priorizar disponibilidad) ante ese escenario de falla. Abadi
extiende esta lectura con el modelo PACELC: incluso \emph{sin} partici\'on,
un sistema distribuido debe resolver un compromiso entre latencia y
consistencia en su operaci\'on normal (\emph{Else}), de modo que un sistema
se clasifica completamente como PC/EC, PA/EL, o combinaciones intermedias
seg\'un su configuraci\'on \cite{abadi2012}.

\subsection{Estrategias de replicaci\'on y consenso}

Las estrategias de replicaci\'on se distinguen por qui\'en acepta las
escrituras. En la replicaci\'on \textbf{primario-secundario} (o
\emph{leader-follower}), una \'unica r\'eplica designada acepta escrituras
y las propaga a las dem\'as, que solo sirven lecturas; ofrece consistencia
fuerte natural pero convierte al primario en un cuello de botella. En la
replicaci\'on \textbf{multi-primario}, varias r\'eplicas aceptan escrituras
concurrentes, mejorando la disponibilidad de escritura a costa de requerir
pol\'iticas de resoluci\'on de conflictos. En la replicaci\'on
\textbf{sin l\'ider} (\emph{leaderless}) basada en \emph{quorum}, un
cliente escribe en paralelo sobre varios nodos y la operaci\'on se
considera exitosa cuando W r\'eplicas confirman, garantiz\'andose lecturas
consistentes si W + R > N \cite{tanenbaum2017}. Los algoritmos de consenso
modernos, en particular Raft, sustituyen a los protocolos de bloqueo como
el commit en dos fases al replicar un registro de decisiones por mayor\'ia
sin asumir sincron\'ia de red, y son la base de sistemas de producci\'on
como etcd y CockroachDB \cite{ongaro2014}.

\subsection{Sistemas de referencia: los dos extremos del espectro}

Dynamo, el almacenamiento clave-valor de Amazon, prioriza la disponibilidad
mediante replicaci\'on sin l\'ider, \emph{quorum} configurable y
consistencia eventual con resoluci\'on de conflictos por vectores de
versi\'on, aceptando que distintas r\'eplicas puedan divergir
temporalmente a cambio de nunca rechazar una escritura \cite{decandia2007}.
En el extremo opuesto, Google Spanner ofrece consistencia externa
(transacciones globalmente serializables) apoy\'andose en TrueTime, una API
de tiempo con incertidumbre acotada mediante relojes at\'omicos y GPS, lo
que le permite comportarse como CP/EC incluso a escala planetaria
\cite{corbett2013}. CockroachDB, la base de datos utilizada en SCLI, se
sit\'ua conceptualmente junto a Spanner: es una base de datos relacional
distribuida que fragmenta autom\'aticamente los datos en \emph{ranges}
replicados por consenso Raft, con aislamiento serializable estricto por
defecto, clasificada como PC/EC bajo PACELC.

% ============================================================
% ANALISIS DEL PFC (Seccion 5) - NUCLEO D2
% ============================================================
\section{An\'alisis del Modelo de Consistencia de SCLI}
\label{sec:analisis}

\subsection{Arquitectura de datos: base de datos por servicio}

SCLI implementa el patr\'on \emph{database per service}: cada
microservicio --- \texttt{auth-service}, \texttt{usuarios-service} y
\texttt{academico-laboratorios-service} (este \'ultimo desarrollado por el
autor de este informe) --- posee su propia instancia CockroachDB aislada,
sin acceso directo a los datos de los dem\'as. La comunicaci\'on entre
servicios ocurre exclusivamente a trav\'es de REST interno protegido con
el encabezado \texttt{X-Internal-Api-Key}. La Figura~\ref{fig:topologia}
ilustra esta topolog\'ia, incluyendo el cuarto servicio,
\texttt{reservas-solicitudes-service}, a\'un incompleto.

\begin{figure}[h]
\centering
\begin{tikzpicture}[
    node distance=0.9cm and 1.1cm,
    service/.style={rectangle, draw, rounded corners, minimum width=2.9cm,
        minimum height=1cm, align=center, fill=blue!8, font=\scriptsize},
    servicepend/.style={rectangle, draw, dashed, rounded corners,
        minimum width=2.9cm, minimum height=1cm, align=center,
        fill=gray!10, font=\scriptsize},
    db/.style={cylinder, draw, shape border rotate=90, aspect=0.25,
        minimum width=1.7cm, minimum height=0.9cm, align=center,
        fill=orange!12, font=\tiny},
    front/.style={rectangle, draw, rounded corners, minimum width=2.9cm,
        minimum height=0.9cm, align=center, fill=green!10, font=\scriptsize}
]

\node[front] (front) {SCLI Front\\(Vue 3)};

\node[service, below=of front, xshift=-4.6cm] (auth) {auth-service\\:8081};
\node[service, below=of front, xshift=-1.5cm] (usu) {usuarios-service\\:8082};
\node[service, below=of front, xshift=1.5cm] (aca) {academico-labs-service\\:8083};
\node[servicepend, below=of front, xshift=4.6cm] (res) {reservas-solicitudes\\:8084 (incompleto)};

\node[db, below=0.55cm of auth] (dbauth) {CRDB\\v26.2.0\\Raft};
\node[db, below=0.55cm of usu] (dbusu) {CRDB\\v24.3.5 $\ast$\\Raft};
\node[db, below=0.55cm of aca] (dbaca) {CRDB\\v26.2.0\\Raft};

\draw[-{Latex[length=2mm]}] (front) -- (auth);
\draw[-{Latex[length=2mm]}] (front) -- (usu);
\draw[-{Latex[length=2mm]}] (front) -- (aca);
\draw[-{Latex[length=2mm]}, dashed] (front) -- (res);

\draw[-{Latex[length=2mm]}] (auth) -- (dbauth);
\draw[-{Latex[length=2mm]}] (usu) -- (dbusu);
\draw[-{Latex[length=2mm]}] (aca) -- (dbaca);

\draw[-{Latex[length=2mm]}, dashed, bend left=15]
    (res.north west) to node[above, font=\tiny] {REST s\'incrono} (usu.east);
\draw[-{Latex[length=2mm]}, dashed, bend right=15]
    (res.north) to node[below, font=\tiny] {REST s\'incrono} (aca.east);

\end{tikzpicture}
\caption{Topolog\'ia de SCLI: consistencia fuerte intra-servicio v\'ia
CockroachDB/Raft (l\'inea continua) frente a coordinaci\'on inter-servicio
sin garant\'ias formales v\'ia REST s\'incrono (l\'inea discontinua).
$\ast$ La versi\'on divergente de CockroachDB en \texttt{usuarios-service}
es un riesgo real de heterogeneidad documentado en el repositorio del
equipo.}
\label{fig:topologia}
\end{figure}

\subsection{Capa intra-servicio: consistencia fuerte heredada de CockroachDB}

El modelo de consistencia \emph{dentro} de cada servicio no fue dise\~nado
expl\'icitamente por el equipo, sino heredado de la elecci\'on tecnol\'ogica
de CockroachDB. Cada base de datos fragmenta autom\'aticamente sus tablas
en \emph{ranges} de aproximadamente 64~MB, cada uno replicado mediante
consenso Raft entre los nodos del cl\'uster, con aislamiento serializable
estricto por defecto. Esto ubica a cada instancia CockroachDB de SCLI como
PC/EC bajo PACELC: ante una partici\'on, el sistema prioriza no devolver
datos potencialmente obsoletos (C sobre A); en operaci\'on normal, el
consenso s\'incrono impone mayor latencia de escritura a cambio de
consistencia estricta (C sobre L) \cite{abadi2012}.

Un hallazgo t\'ecnico propio refuerza esta caracterizaci\'on: la
integraci\'on de Flyway con CockroachDB v26.2 en
\texttt{academico-laboratorios-service} requiri\'o agregar el par\'ametro
JDBC \texttt{options=-c\%20allow\_unsafe\_internals\%3Dtrue}, evidencia
directa de que los cambios de esquema deben coordinarse con las garant\'ias
internas del motor y no pueden tratarse como operaciones triviales en un
sistema con consenso distribuido. La \textbf{divergencia de versi\'on}
entre \texttt{usuarios-service} (CockroachDB v24.3.5) y el resto de
servicios (v26.2.0), visible en la Figura~\ref{fig:topologia}, es adem\'as
un riesgo concreto: aunque cada base opera de forma aislada y esto no
compromete la consistencia \emph{interna} de ninguna de ellas, s\'i
introduce heterogeneidad de comportamiento (posibles diferencias en
optimizador de consultas o semántica de \texttt{ALTER TABLE}) que
complicar\'ia una futura federaci\'on de esquemas entre servicios.

\subsection{Capa inter-servicio: coordinaci\'on sin garant\'ias formales}

A diferencia de la capa intra-servicio, la comunicaci\'on \emph{entre}
microservicios de SCLI no implementa ning\'un modelo de consistencia
expl\'icito. Los clientes REST de \texttt{reservas-solicitudes-service}
(\texttt{AcademicoLaboratoriosClient}, \texttt{UsuariosClient}, construidos
sobre \texttt{RestClient} de Spring) realizan llamadas s\'incronas de
solicitud-respuesta sin coordinaci\'on transaccional, sin patr\'on de
compensaci\'on y sin cola de mensajes intermedia. El \'unico mecanismo de
consistencia a nivel de aplicaci\'on documentado por el equipo es
\textbf{local} a \texttt{reservas-solicitudes-service}: concurrencia
optimista mediante un campo \texttt{version} sobre las entidades
\texttt{SolicitudReserva} y \texttt{Reserva}, e idempotencia en las
operaciones de iniciar, finalizar y cancelar una reserva, seg\'un el
documento de dise\~no \texttt{02-maquina-estados.md} del propio equipo. La
coordinaci\'on entre la aprobaci\'on de una solicitud y la creaci\'on de su
reserva correspondiente se resuelve \'unicamente dentro de una transacci\'on
local a ese servicio; no existe ning\'un mecanismo an\'alogo para
coordinar, por ejemplo, una eventual inconsistencia si
\texttt{academico-laboratorios-service} cambia el estado de un laboratorio
mientras \texttt{reservas-solicitudes-service} procesa una reserva sobre
\'el.

En t\'erminos de PACELC, esta capa no puede clasificarse formalmente
porque no fue dise\~nada como una decisi\'on consciente: es, de facto, un
acoplamiento s\'incrono fr\'agil sin las garant\'ias de un modelo PC (no
hay consenso ni transacci\'on distribuida) ni la resiliencia expl\'icita de
un modelo AP (no hay degradaci\'on controlada ni cola de reintento ante
fallo de un servicio dependiente).

% ============================================================
% EVALUACION DE ALTERNATIVAS (Seccion 6) - NUCLEO D3
% ============================================================
\section{Evaluaci\'on de Alternativas}

Se contrasta el modelo actual (REST s\'incrono sin coordinaci\'on expl\'icita
entre servicios, consistencia fuerte solo intra-servicio) con dos
alternativas de dise\~no aplicables al dominio de SCLI.

\textbf{Alternativa A --- Propagaci\'on por eventos (\emph{outbox} +
consistencia eventual inter-servicio).} Cada servicio publica eventos de
dominio (por ejemplo, \texttt{LaboratorioActualizado},
\texttt{PerfilModificado}) mediante el patr\'on \emph{transactional
outbox}, consumidos as\'incronamente por los servicios interesados para
mantener copias locales de solo lectura. La consistencia intra-servicio se
mantiene fuerte (CockroachDB/Raft); la consistencia inter-servicio pasa a
ser expl\'icitamente eventual.

\textbf{Alternativa B --- Orquestaci\'on tipo Saga con compensaci\'on.}
Para operaciones cr\'iticas que cruzan servicios --- como la aprobaci\'on
de una solicitud que depende de validar disponibilidad en
\texttt{academico-laboratorios-service} y el rol del usuario en
\texttt{usuarios-service} --- se introduce un orquestador de Saga que
ejecuta pasos remotos y, ante un fallo parcial, dispara acciones
compensatorias expl\'icitas en lugar de asumir at\'omicidad inexistente.

La Tabla~\ref{tab:alternativas} compara ambas alternativas con el modelo
actual seg\'un cinco criterios t\'ecnicos.

\begin{table}[h]
\centering
\caption{Comparaci\'on del modelo actual frente a dos alternativas de
consistencia y replicaci\'on inter-servicio}
\label{tab:alternativas}
\small
\begin{tabular}{@{}p{3.1cm}p{3.1cm}p{3.1cm}p{3.1cm}@{}}
\toprule
\textbf{Criterio} & \textbf{Actual (REST s\'incrono)} & \textbf{A: Eventos} & \textbf{B: Saga} \\
\midrule
Consistencia inter-servicio & Indefinida & Eventual & Fuerte por operaci\'on \\
Disponibilidad ante fallo remoto & Baja (falla en cascada) & Alta & Media \\
Latencia percibida & Baja (si todo responde) & Muy baja (lectura local) & Media--alta \\
Tolerancia a partici\'on & No dise\~nada & Alta & Media \\
Complejidad de implementaci\'on & Baja & Media & Alta \\
\bottomrule
\end{tabular}
\end{table}

\subsection{Postura cr\'itica: \textquestiondown es el modelo actual \'optimo?}

La respuesta es \textbf{parcialmente s\'i, parcialmente no}, y depende de
separar las dos capas identificadas en la Secci\'on~\ref{sec:analisis}.

\textbf{La capa intra-servicio s\'i es una elecci\'on \'optima o cercana a
\'optima} para el dominio de laboratorios inform\'aticos. Datos como el
estado de un laboratorio o la asignaci\'on de un equipo requieren evitar
condiciones de sobre-asignaci\'on --- an\'alogas al riesgo de sobregiro en
un dominio bancario --- por lo que la consistencia fuerte que CockroachDB
ofrece por defecto, v\'ia Raft y aislamiento serializable, est\'a
justificada por el dominio y no por inercia tecnol\'ogica.

\textbf{La capa inter-servicio, en cambio, no es \'optima en su estado
actual}, no porque el REST s\'incrono sea inv\'alido en general, sino
porque el equipo no tom\'o una decisi\'on consciente sobre qu\'e modelo de
consistencia aplicar entre servicios: el acoplamiento s\'incrono actual
combina las desventajas de un modelo fuerte (latencia acumulada, fallo en
cascada) sin ofrecer ninguna de sus garant\'ias reales (no hay
at\'omicidad cross-service). Esta limitaci\'on se agrava porque
\texttt{reservas-solicitudes-service} --- el componente que m\'as depende
de esta coordinaci\'on --- ni siquiera compila actualmente, al faltar las
implementaciones de \texttt{ReservaService} y
\texttt{SolicitudReservaService}; el equipo est\'a a tiempo de decidir el
modelo de coordinaci\'on \emph{antes} de terminar de conectar esas piezas,
en lugar de despu\'es.

Como mitigaci\'on concreta, se recomienda adoptar la \textbf{Alternativa
B (Saga con compensaci\'on)} espec\'ificamente para el flujo cr\'itico de
aprobaci\'on de solicitudes --- donde una inconsistencia s\'i tendr\'ia
impacto operativo real --- y reservar la \textbf{Alternativa A (eventos)}
para la propagaci\'on de datos de solo lectura de bajo riesgo, como copias
locales del perfil de usuario o del estado general de un laboratorio, que
toleran consistencia eventual sin afectar la integridad del dominio.

% ============================================================
% CONCLUSIONES (Seccion 7)
% ============================================================
\section{Conclusiones}

El modelo de consistencia de SCLI no es una decisi\'on \'unica sino dos
decisiones superpuestas: una fuerte y bien justificada, heredada de
CockroachDB dentro de cada microservicio; y otra, impl\'icita y no
dise\~nada, en la coordinaci\'on entre servicios v\'ia REST s\'incrono. El
teorema CAP y el modelo PACELC permiten clasificar con precisi\'on la
primera capa (PC/EC) pero exponen la ausencia de una clasificaci\'on
equivalente para la segunda, que es precisamente el hallazgo cr\'itico de
este an\'alisis. Se recomienda que el equipo adopte expl\'icitamente un
patr\'on de coordinaci\'on --- Saga para operaciones cr\'iticas, eventos
para propagaci\'on de solo lectura --- antes de completar la
implementaci\'on de \texttt{reservas-solicitudes-service}, transformando
una limitaci\'on arquitect\'onica actual en una decisi\'on de dise\~no
consciente y defendible.

% ============================================================
% DECLARACION DE USO DE IA (Seccion 8)
% ============================================================
\section*{Declaraci\'on de Uso de Inteligencia Artificial}

Se utiliz\'o el asistente de inteligencia artificial Claude (Anthropic)
como apoyo en: (i) la estructuraci\'on del documento conforme a los
requisitos del TA-IND-02; (ii) la redacci\'on y organizaci\'on de las
secciones de marco te\'orico, an\'alisis y evaluaci\'on de alternativas a
partir de la informaci\'on real del repositorio del PFC proporcionada por
el autor; (iii) la elaboraci\'on del diagrama de topolog\'ia (Figura~1) en
c\'odigo TikZ. La caracterizaci\'on t\'ecnica del sistema SCLI, la
verificaci\'on de las referencias citadas y la postura cr\'itica final
fueron revisadas y validadas por el autor.

% ============================================================
% REFERENCIAS (Seccion 9)
% ============================================================
\begin{thebibliography}{00}

\bibitem{abadi2012} D. J. Abadi, ``Consistency tradeoffs in modern
distributed database system design: CAP is only part of the story,''
\emph{IEEE Computer}, vol. 45, no. 2, pp. 37--42, Feb. 2012, doi:
10.1109/MC.2012.33.

\bibitem{corbett2013} J. C. Corbett et al., ``Spanner: Google's globally
distributed database,'' \emph{ACM Transactions on Computer Systems},
vol. 31, no. 3, art. no. 8, pp. 1--22, 2013, doi: 10.1145/2491245.

\bibitem{decandia2007} G. DeCandia et al., ``Dynamo: Amazon's highly
available key-value store,'' in \emph{Proc. 21st ACM SIGOPS Symp.
Operating Systems Principles (SOSP '07)}, 2007, pp. 205--220, doi:
10.1145/1294261.1294281.

\bibitem{gilbert2002} S. Gilbert and N. Lynch, ``Brewer's conjecture and
the feasibility of consistent, available, partition-tolerant web
services,'' \emph{ACM SIGACT News}, vol. 33, no. 2, pp. 51--59, Jun. 2002,
doi: 10.1145/564585.564601.

\bibitem{viotti2016} P. Viotti and M. Vukoli\'c, ``Consistency in
non-transactional distributed storage systems,'' \emph{ACM Computing
Surveys}, vol. 49, no. 1, art. no. 19, pp. 1--34, 2016, doi:
10.1145/2926965.

\bibitem{ongaro2014} D. Ongaro and J. Ousterhout, ``In search of an
understandable consensus algorithm,'' in \emph{Proc. 2014 USENIX Annual
Technical Conference (USENIX ATC '14)}, Philadelphia, PA, USA, Jun. 2014,
pp. 305--319, isbn: 978-1-931971-10-2.

\bibitem{tanenbaum2017} M. van Steen and A. S. Tanenbaum,
\emph{Distributed Systems}, 3rd ed. CreateSpace Independent Publishing
Platform, 2017, isbn: 978-1543057386.

\bibitem{kleppmann2017} M. Kleppmann, \emph{Designing Data-Intensive
Applications: The Big Ideas Behind Reliable, Scalable, and Maintainable
Systems}, 1st ed. Sebastopol, CA, USA: O'Reilly Media, 2017, isbn:
978-1-4493-7332-0.

\end{thebibliography}

\end{document}